% Generated from paper/full_draft. Edit the Markdown source or migration script.

\section{Limitations and Future Work}
\label{sec:08-limitations:limitations-and-future-work}

\subsection{Simulated Feedback}
\label{sec:08-limitations:simulated-feedback}

The experiments derive simulated feedback from ground truth under controlled noise and trust. They test policy response to a feedback signal, not behavior under delayed, biased, adversarial, low-quality, or non-stationary human input. The GT-derived hard-case experiment is likewise a recovery test after a failed compacted answer, not first-pass IID improvement. It repairs some cases whose evidence remains reachable through the candidate pool and arms, without implying universal recovery.

After training route state on disjoint history folds and freezing it, learned full routing does not outperform matched static-nearest or cold no-feedback full routing on unseen queries in either LoTTE domain; learned gating is significantly below Dense in all three seeds of both. Feedback is therefore a controlled repeated-interaction and recovery mechanism, not a demonstrated universal first-pass gain.

\subsection{Limited Generation Evaluation}
\label{sec:08-limitations:limited-generation-evaluation}

Across 300 frozen-test queries, seven methods, 2,100 answers, and 6,272 valid three-judge ratings, matched BGE, E5, and SentMMR comparisons save context without a statistically detectable correctness change. However, one model generates all answers, DeepSeek is also a judge, no human ratings are available, and only one LoTTE domain is covered. MiniMax-M3 content filtering leaves 28 judgments unavailable, which cross-judge analyses exclude rather than impute. Calibration differs across judges, and majority faithfulness falls for BGE but rises for SentMMR. The evidence supports bounded correctness robustness, not strict non-inferiority, uniform faithfulness, answer superiority, or user satisfaction.

\subsection{Dense Baseline, Evidence Completeness, and Domain Calibration}
\label{sec:08-limitations:dense-baseline-evidence-completeness-and-domain-calibration}

Dense-only retrieval remains a strong baseline and recall floor, not a route that IntentRoute universally replaces. Moreover, preserving query-level $\mathrm{Hit@10}$ does not guarantee coverage of all ground-truth chunks. Legal review, medical synthesis, and other complete-evidence tasks may require a conservative policy or no compaction. The supported result is therefore a calibrated quality-context frontier, not universal Dense replacement or lossless evidence compression.

LoTTE science/search transfers fixed top-10 ranking gains, but not budget strength: at 100k, an aggressive budget saves 17--21\% while introducing small frozen-test $\mathrm{Hit@10}$ drops. Budgets require domain- and scale-specific calibration with dense fallback for low-confidence or high-risk regions.

Across 20 deterministic partitions, selected-policy mean Hit remains within \nolinkurl{1pp} of dense on every 200k/638k test partition, but only 70\%/85\% of 100k/400k partitions. Because partitions overlap, they measure sensitivity rather than independent replication. The original frozen split remains valid, but deployment should prefer repeated or nested calibration over a universal no-loss inference.

The disjoint five-fold follow-up yields 14.50\% mean 400k saving at effectively zero mean Hit delta, yet fold deltas remain heterogeneous and every fold selects a different policy. It addresses missing calibration without establishing split-invariant behavior.

The prospectively specified 100k expansion confirms domain effects. No-feedback writing/search yields 10.09\% saving, +0.12pp mean Hit, and 2/3 strict non-inferiority seeds; recreation/search yields 5.42\%, -0.76pp, and 0/3. Trust-weighted calibration falls back to Dense in every fold of both. These are domain-specific operating points, not a universal quality-preserving token-saving guarantee.

Seeds 13/17/19 are engineering replicates; query-level paired tests provide the main inference. The 400k saving interval remains wider than at other scales, and cross-fitting establishes strict non-inferiority in 0/3 seeds.

The fixed \nolinkurl{1pp} threshold is an engineering guardrail, not an equivalence theorem: it represents about four original-split or six out-of-fold hit events. Scales, controls, backbones, datasets, and judges map heterogeneous mechanisms rather than form an IID replication pool, so secondary $p$-values are neither aggregated nor used for global superiority.

\subsection{Geometry and Fixed-Arm Scope}
\label{sec:08-limitations:geometry-and-fixed-arm-scope}

Nearest-cluster hit, PCA spectrum, and context retention support informative local routing geometry, not a manifold theorem. Dense retrieval remains necessary, and strong cluster-local signals in recreation/search and writing/search coexist with different budget outcomes; geometry is not a direct compression-safety predictor.

KMeans/MiniBatchKMeans supplies the fixed, reproducible, scalable arm space required by the tested LinUCB setup, not a universally optimal clustering method. HDBSCAN or graph clusters may suit deployments but introduce dynamic arm management. Across $K=8$--$128$, full multi-route quality is stable and gating is sensitive, making $K=32$ an engineering point rather than an optimum.

\subsection{Limited Baseline, Encoder, and Domain Coverage}
\label{sec:08-limitations:limited-baseline-encoder-and-domain-coverage}

Matched MiniLM, BGE-base, and E5-base backbones and a cross-encoder support robustness within LoTTE, but exclude domain-specific, late-interaction, and proprietary encoders. Only BGE has an above-dense quality-first point; E5 supplies a near-dense saving point. Sentence-MMR is the shared downstream compressor, and the official LLMLingua-2 follow-up supports the same Dense-versus-IntentRoute direction on one frozen 300-query setting. Its three-judge correctness intervals cross zero, so it is compressor-complementarity evidence rather than strict non-inferiority, human validation, or broad parity with learned token-level prompt compression.

Of nine settings, only LoTTE technology/search receives the complete multi-scale, matched-baseline, calibration/test, and generation protocol. Science/search adds domain and scale transfer; recreation/search and writing/search add common-protocol 100k retrieval/budget rows without multi-scale or answer evaluation; biomedical datasets, Banking77, eManual, and CUAD test transfer, mechanism, or boundaries. They are not nine independent full-stack replications, and more repeated vertical-corpus evaluations remain necessary.

\subsection{Future Work}
\label{sec:08-limitations:future-work}

The attribution boundary and class imbalance in Section 6.5 motivate, but do not validate, a learned per-query confidence-to-token-ratio predictor.

Future work should evaluate:

\begin{itemize}
\item real user feedback with trust scoring and delayed-feedback handling;
\item multi-model and human-rated answer-quality and citation-faithfulness studies;
\item stronger dense encoders, rerankers, and late-interaction retrieval models;
\item additional domains, generators, and human raters for learned-compressor evaluation;
\item graph or density-based dynamic clustering under bandit-compatible arm management;
\item repeated evaluations on additional vertical-domain corpora;
\item production policies that lower dense usage only after route confidence is demonstrably stable;
\item recovery policies evaluated with real delayed feedback rather than GT-derived same-query retry.
\end{itemize}
